\itemhead{复习题三 · 第12题（教材p.91）}{XJ-B1-C03-RV-P091-Y0650}
\label{XJ-B1-C03-RV-P091-Y0650}
{\footnotesize 湘教版教材p.91 / PDF第98页 / 复习题三 / 第12题}
\field{原题}{（多选题）已知函数 $f(x)$ 满足 $f(x+y)=f(x)+f(y)\ (x,y\in\mathbb R)$，则（\quad）
\choice{A}$f(0)=0$
\choice{B}$f(x)$ 是奇函数
\choice{C}$f(x)$ 是偶函数
\choice{D}当 $f(2)=\sqrt5$ 时，$f(200)=100\sqrt5$}
\field{参考答案}{ABD。令 $x=y=0$ 得 $f(0)=0$；令 $y=-x$ 得奇性；$f(200)=100f(2)=100\sqrt5$。C不必成立，$f(x)=x$ 为反例。}
\field{核对说明}{多选设问变化：现教材需写ABD，答案册PDF25按三个证明小问给解，须转换为逐项判断。}
{\small \field{关联知识}{F08 奇偶性、对称与函数性质应用（教材帮p.113、114、115 / PDF122、123、124）；F03 函数求值、解析式与函数建模（教材帮p.80、87、88、89、90、131 / PDF89、96、97、98、99、140）}
\field{方法与易错点}{检查定义域，再代入 $-x$；求另一侧解析式时同时写出 $-x$ 所在区间。 易错：只看公式不看定义域；用一对点证明奇偶；负半轴代换遗漏最外层负号。}
\field{答案依据}{答案册相关题组：PDF第24、25、26、27页（册内页23、24、25、26）；按当前题库原题逐项重算，不按旧版题号机械对套。}}
